\chapter{可靠性门控的伪标签学习}

\section{模块定位}
第三章解决了目标域样本如何通过教师--学生一致性进入训练的问题，但目标域伪标签本身仍可能存在噪声。对于 29 类流程工业故障诊断任务，简单采用最大概率作为伪标签采纳依据，容易把高置信错误样本引入训练，并进一步造成错误监督的累积。与此同时，不同故障类别在目标域上的预测置信度并不均衡，若采用单一全局阈值，伪标签预算很可能被少数容易类别占据。

本章对应第一章提出的挑战二：目标域伪标签容易受噪声和类别不均衡影响。为此，RCTA 在教师--学生一致性的基础上进一步引入多因素可靠性评分、按类门控和分层利用机制，使目标域伪监督从“看到样本就直接使用”转变为“对样本质量分层评估、分层利用”。

\section{伪标签误差传播分析}
在流程工业多类别诊断中，伪标签误差传播通常来自以下三个来源。

\begin{enumerate}
\item \textbf{过度自信错误。} 某些目标样本虽然预测最大概率较高，但实质上位于错误类别簇附近，此时单纯以置信度作为采纳依据会把错误监督以高权重注入训练。
\item \textbf{边界样本不稳定。} 对于靠近类别边界或跨域偏移较大的样本，弱增强和强增强视图下的预测可能并不一致，若忽略这种不稳定性，伪监督容易放大边界噪声。
\item \textbf{类别预算失衡。} 在 29 类任务中，一些容易类别会较早形成较高置信度，若采用全局 top-k 或单一阈值，少数类别可能占满伪标签预算，导致困难类别长期得不到纠偏机会。
\end{enumerate}

因此，可靠性门控模块的核心目标不是简单“筛掉低分样本”，而是在训练不同阶段尽可能用有限的目标域伪监督预算换取更高质量的监督信号，并保持类别层面的覆盖性。

\section{多因素可靠性评分}
目标域伪标签是否可用，不能只由单一最大概率决定。为此，RCTA 对每个目标样本构造四个互补的可靠性分量：
\begin{align}
q_j &= \max p_j^{tea},\\
e_j &= 1-\frac{-\sum_{k=1}^{K}p_{jk}^{tea}\log p_{jk}^{tea}}{\log K},\\
u_j &= p_{j,\hat{y}_j^t}^{stu},\\
s_j &= \frac{1+\cos(f_j^{tea},c_{\hat{y}_j^t}^{ref})}{2}.
\end{align}
其中，$q_j$ 表示教师置信度，$e_j$ 表示逆归一化熵，$u_j$ 表示学生在强增强视图上对伪标签类别的响应强度，$s_j$ 表示教师特征与参考原型之间的相似度。于是综合可靠性得分定义为
\begin{equation}
r_j=
\omega_1 q_j+\omega_2 e_j+\omega_3 u_j+\omega_4 s_j,
\qquad
\sum_{m=1}^{4}\omega_m=1.
\end{equation}

在当前 fixed-fold benchmark 配置中，四个分量的原始权重依次设置为 1.1、1.0、1.1 和 1.2，并在实现中归一化后参与加权。计算完成后，还会进一步施加幂次变换
\begin{equation}
\tilde{r}_j=r_j^{\gamma},
\end{equation}
其中当前配置中的 $\gamma=1.4$，目的是适度放大高可靠样本与中低可靠样本之间的区分度。

这样的设计具有明显互补性。若只依赖最大概率，模型容易把“高置信但错误”的样本误认为可靠；若只依赖熵，难以区分边界犹豫样本与真实混乱样本；若只依赖增强一致性，又会忽略类别结构位置；若只依赖原型相似度，则在原型尚未稳定时容易引入偏差。将四种信息联合起来，可以在“预测尖锐度、视图稳定性、结构匹配度”三个层面同时评估样本质量。

\section{按类门控与样本接纳课程}
在获得可靠性得分之后，RCTA 并不让所有目标域样本都参与伪标签监督，而是同时引入分数下界和按类接纳比例两个门控量。分数下界定义为
\begin{equation}
\tau(t)=\tau_{start}
+(\tau_{end}-\tau_{start})\cdot
\min\left(\frac{t}{T_{\tau}},1\right),
\end{equation}
按类接纳比例定义为
\begin{equation}
\rho(t)=\rho_{start}
+(\rho_{end}-\rho_{start})\cdot
\min\left(\frac{t}{T_{\rho}},1\right).
\end{equation}
在当前 benchmark 配置中，$\tau(t)$ 从 0.34 线性上升到 0.56，$\rho(t)$ 从 0.22 线性上升到 0.72，对应的课程步数均为 2200。其核心意图是：训练早期只保留少量高可靠样本，避免噪声伪监督主导训练；训练中后期再逐步扩大接纳范围，使模型有机会利用更多目标域信息。

与简单的全局阈值不同，RCTA 采用按预测类别分别筛选的策略。对每个预测类别 $k$，先保留满足 $\tilde{r}_j\ge \tau(t)$ 的样本，再在该类别内部选取得分前 $\lceil \rho(t)\cdot n_k\rceil$ 个样本进入门控集合 $\Omega_t$。这种按类门控方式能够缓解类别不均衡带来的“头部类别占满伪标签预算”问题，更适合 29 类流程工业诊断任务。

\section{分层利用：可靠、半可靠与不可靠样本}
当前实现并不把目标域样本简单分为“可用”和“不可用”两类，而是进一步划分为可靠、半可靠和不可靠三组。设可靠阈值和半可靠阈值分别为 $\eta_r$ 与 $\eta_s$，则有
\begin{align}
\Omega_t^{r} &= \{j\mid j\in\Omega_t,\ \tilde{r}_j\ge \eta_r\},\\
\Omega_t^{s} &= \{j\mid j\in\Omega_t,\ \eta_s\le \tilde{r}_j < \eta_r\},\\
\Omega_t^{u} &= \{1,\ldots,m_t\}\setminus(\Omega_t^{r}\cup\Omega_t^{s}).
\end{align}
在当前 benchmark 配置中，$\eta_r=0.66$，$\eta_s=0.44$。这种分层利用机制的意义在于：同样被 gate 接纳的样本，其可靠程度仍然不同，完全使用同一强度监督并不合理。

具体而言，可靠样本主要承担伪标签监督与较强的一致性约束；半可靠样本主要承担较弱的一致性正则；不可靠样本则不进入伪标签交叉熵，而是通过轻量熵项维持适度结构化。这样做的好处在于，模型不会因为“分数刚刚过线”的样本而获得与高可靠样本同等强度的监督，也不会彻底浪费那些尚未达到高可靠标准、但已经包含一定信息的样本。

\section{可靠性加权伪标签损失与辅助正则}
对于通过门控的目标样本，伪标签损失采用带可靠性加权的交叉熵~\cite{lee2013pseudolabel}：
\begin{equation}
\mathcal{L}_{pl}=
\frac{1}{\sum_{j\in\Omega_t}\tilde{r}_j}
\sum_{j\in\Omega_t}
\tilde{r}_j\,
\ell_{ce}\!\left(h(\phi(\tilde{x}_{j,s}^{t})),\hat{y}_j^t\right).
\end{equation}
这样一来，不同目标样本虽然都被允许进入训练，但高可靠样本会对梯度更新产生更大影响。

除主伪标签损失外，当前实现还额外保留两项与可靠性分层相关的辅助正则。对半可靠样本，施加较弱的一致性项：
\begin{equation}
\mathcal{L}_{semi-cons}=
\frac{1}{|\Omega_t^{s}|}\sum_{j\in\Omega_t^{s}}
\tilde{r}_j\,
\left\|p_j^{stu}-p_j^{tea}\right\|_2^2.
\end{equation}
对不可靠样本，则使用轻量熵项维持预测结构化：
\begin{equation}
\mathcal{L}_{u-ent}=
\frac{1}{|\Omega_t^{u}|}\sum_{j\in\Omega_t^{u}}
\left(
-\sum_{k=1}^{K}p_{jk}^{stu}\log p_{jk}^{stu}
\right).
\end{equation}
在当前 benchmark 配置中，这两项对应的附加权重分别为 0.04 和 0.01。它们的目标不是单独主导训练，而是对“还不够可靠、但也不是完全无用”的样本提供更温和的处理方式。

\section{与前后模块的耦合关系}
可靠性门控位于第三章和第五章之间，既依赖前序模块，也会反向影响后序模块。首先，本章的教师概率、学生强增强响应和参考原型相似度均建立在第三章输出的稳定概率与特征之上；没有第三章的一致性底座，可靠性得分本身会更不稳定。其次，本章筛选出的高可靠样本会进入第五章的目标域伪监督与原型约束，而第五章构造的参考原型又会反过来参与本章的结构相似度评分。

从这个角度看，本章的真正贡献并不仅仅是一个门槛函数，而是把目标域样本利用过程从“单点判断”升级为“基于概率、一致性与结构信息的闭环筛选”。这也是 RCTA 能够避免把所有目标域样本一股脑塞进训练的重要原因。

\section{本章小结}
本章围绕可靠性门控的伪标签学习展开，分别说明了伪标签误差传播的来源、多因素可靠性评分、按类门控课程、可靠/半可靠/不可靠三层样本划分，以及可靠性加权伪标签损失与辅助正则。该模块的核心价值在于提高目标域伪监督的质量和利用效率，并为第五章的类条件结构约束提供更干净、更可控的目标域样本来源。
